\documentclass[12pt,a4paper]{article}

% ============================================================
% PACKAGES
% ============================================================

\usepackage[utf8]{inputenc}
\usepackage[T1]{fontenc}
\usepackage[french]{babel}

\usepackage[a4paper,
    top=1.7cm,
    bottom=1.7cm,
    left=2cm,
    right=2cm]{geometry}

\usepackage{graphicx}
\usepackage{xcolor}
\usepackage{tikz}
\usepackage{setspace}
\usepackage{lmodern}
\usepackage{microtype}
\usepackage{amsmath}
\usepackage{float}
\usepackage[most]{tcolorbox}
\usepackage{booktabs}

% ============================================================
% COULEURS
% ============================================================

\definecolor{BleuPrincipal}{RGB}{20,55,95}
\definecolor{BleuClair}{RGB}{225,235,245}
\definecolor{GrisTexte}{RGB}{70,70,70}

% ============================================================
% INFORMATIONS À MODIFIER
% ============================================================

\newcommand{\universite}
{Université Cheikh Anta Diop de Dakar}

\newcommand{\faculte}
{Faculté des Sciences et Techniques}

\newcommand{\departement}
{Département de Mathématiques et Informatique}

\newcommand{\laboratoire}
{Laboratoire d'Algèbre de Cryptologie}

\newcommand{\laboratoireCourt}
{LAGCAA}

\newcommand{\master}
{%
Master Transmission de Données et Sécurité de l'Information\\[0.15cm]
{\normalsize\itshape Option : Cryptographie-Informatique}%
}

\newcommand{\titre}
{%
\textbf{Analyse comparative des modes d'opération d'AES}\\[0.15cm]
\textbf{(CBC vs GCM)}\\[0.15cm]
Sécurité face aux attaques actives et performances
dans la transmission de données%
}

\newcommand{\auteur}
{Koffi Cyril SEPLEDJI }

\newcommand{\encadreur}
{Dr. Babacar Alassane NDAW}

\newcommand{\dateSoutenance}
{Le 17 Mai 2026}

\newcommand{\annee}
{Année académique 2025--2026}

% ============================================================
% LOGOS
% ============================================================

% Si les images existent, elles sont affichées.
% Sinon, un espace réservé est affiché.

\newcommand{\logoUniversite}{logo-universite.png}
\newcommand{\logoLaboratoire}{logo-laboratoire.png}

\newcommand{\afficherLogoUniversite}{%
    \IfFileExists{\logoUniversite}{%
        \includegraphics[
            width=2.5cm,
            height=2.5cm,
            keepaspectratio
        ]{\logoUniversite}%
    }{%
        \begin{tikzpicture}
            \draw[gray] (0,0) rectangle (2.5,2.5);
            \node[align=center,text width=2.2cm]
            at (1.25,1.25)
            {\scriptsize Logo\\Université};
        \end{tikzpicture}%
    }%
}

\newcommand{\afficherLogoLaboratoire}{%
    \IfFileExists{\logoLaboratoire}{%
        \includegraphics[
            width=2.5cm,
            height=2.5cm,
            keepaspectratio
        ]{\logoLaboratoire}%
    }{%
        \begin{tikzpicture}
            \draw[gray] (0,0) rectangle (2.5,2.5);
            \node[align=center,text width=2.2cm]
            at (1.25,1.25)
            {\scriptsize Logo\\Laboratoire};
        \end{tikzpicture}%
    }%
}

% ============================================================
% DOCUMENT
% ============================================================

\begin{document}

\thispagestyle{empty}

% ============================================================
% CADRE DE LA PAGE
% ============================================================

\begin{tikzpicture}[remember picture,overlay]

    % Cadre extérieur
    \draw[
        color=BleuPrincipal,
        line width=1.2pt
    ]
    ([xshift=0.8cm,yshift=-0.8cm]current page.north west)
    rectangle
    ([xshift=-0.8cm,yshift=0.8cm]current page.south east);

    % Cadre intérieur
    \draw[
        color=BleuPrincipal!50,
        line width=0.5pt
    ]
    ([xshift=1.0cm,yshift=-1.0cm]current page.north west)
    rectangle
    ([xshift=-1.0cm,yshift=1.0cm]current page.south east);

\end{tikzpicture}

% ============================================================
% EN-TÊTE
% ============================================================

\begin{center}

\vspace*{0.15cm}

\begin{minipage}[c]{0.20\textwidth}
    \centering
    \afficherLogoUniversite
\end{minipage}
\hfill
\begin{minipage}[c]{0.54\textwidth}
    \centering

    {\large\bfseries
    \textcolor{BleuPrincipal}{\universite}
    \par}

    \vspace{0.12cm}

    {\normalsize\bfseries
    \faculte
    \par}

    \vspace{0.05cm}

    {\small
    \departement
    \par}

\end{minipage}
\hfill
\begin{minipage}[c]{0.20\textwidth}
    \centering
    \afficherLogoLaboratoire
\end{minipage}

\vspace{0.25cm}

% Ligne bleue
{\textcolor{BleuPrincipal}
{\rule{0.88\textwidth}{1.2pt}}}

\vspace{0.15cm}

{\small\bfseries
\textcolor{BleuPrincipal}{\laboratoire}
\par}

{\scriptsize
\laboratoireCourt
\par}

\vspace{0.25cm}

% ============================================================
% MASTER
% ============================================================

{\large\itshape\bfseries
\master
\par}

\vspace{0.2cm}

{\textcolor{BleuPrincipal}
{\rule{0.48\textwidth}{0.7pt}}}

% ============================================================
% SUJET
% ============================================================

\vspace{0.55cm}

{\large\bfseries
Sujet :
\par}

\vspace{0.3cm}

% Boîte du titre

\begin{tikzpicture}

\node[
    draw=BleuPrincipal,
    line width=1pt,
    rounded corners=6pt,
    fill=BleuClair,
    inner xsep=0.55cm,
    inner ysep=0.45cm,
    text width=0.80\textwidth,
    align=center
]
{
    \begin{minipage}{0.76\textwidth}
        \centering
        \setstretch{1.15}

        {\Large\bfseries
        \textcolor{BleuPrincipal}{\titre}
        }

    \end{minipage}
};

\end{tikzpicture}

% ============================================================
% AUTEUR
% ============================================================

\vspace{0.7cm}

{\normalsize
\textit{Projet présenté par}
\par}

\vspace{0.12cm}

{\Large\bfseries
\auteur
\par}

% ============================================================
% ENCADREUR
% ============================================================

\vspace{0.45cm}

{\normalsize
\textit{Sous la direction de}
\par}

\vspace{0.1cm}

{\large\bfseries
\encadreur
\par}

% ============================================================
% BAS DE PAGE
% ============================================================

\vfill

{\textcolor{BleuPrincipal}
{\rule{0.48\textwidth}{0.7pt}}}

\vspace{0.2cm}

{\normalsize\bfseries
\dateSoutenance
\par}

\vspace{0.12cm}

{\small
\textcolor{GrisTexte}{\annee}
\par}

\vspace{0.1cm}

\end{center}


% ==============================
% TABLE DES MATIÈRES
% ==============================
\newpage

\tableofcontents

\newpage
% ==============================
% DÉBUT DU MÉMOIRE
% ==============================

\section{Introduction}

La sécurisation des données échangées dans les systèmes numériques repose largement sur l'utilisation de mécanismes cryptographiques capables d'assurer leur confidentialité et leur intégrité. Parmi les algorithmes de chiffrement symétrique les plus utilisés, l'\textit{Advanced Encryption Standard} (AES) occupe une place importante grâce à son niveau de sécurité, son efficacité et ses nombreuses applications dans les systèmes informatiques et les réseaux.

Cependant, la sécurité d'une communication ne dépend pas uniquement de l'algorithme AES, mais également du mode d'opération utilisé et de la manière dont celui-ci est mis en œuvre. Cette étude s'intéresse particulièrement aux modes CBC (\textit{Cipher Block Chaining}) et GCM (\textit{Galois/Counter Mode}), qui présentent des mécanismes et des propriétés de sécurité différents. Alors que CBC assure principalement le chiffrement des données et nécessite un mécanisme complémentaire pour garantir leur intégrité, GCM associe directement chiffrement et authentification dans un même mécanisme.

Cette différence soulève également des questions concernant la résistance de ces deux modes face aux attaques actives, ainsi que leur influence sur les performances d'une communication sécurisée. Le choix d'un mode d'opération doit en effet prendre en compte à la fois les exigences de sécurité et les contraintes liées au temps de traitement et au volume de données à transmettre.

La problématique étudiée est donc la suivante :

\begin{quote}
\textit{Quel est l'impact du choix entre les modes CBC et GCM sur la sécurité d'une communication face aux attaques actives, et quelles différences de performances peut-on observer lors de la transmission des données ?}
\end{quote}

L'objectif de ce travail est de comparer ces deux modes à travers une étude théorique et expérimentale portant sur leur fonctionnement, leur sécurité face aux attaques actives et leurs performances. Une attention particulière sera accordée aux mécanismes d'authentification, aux principales vulnérabilités associées à chaque mode ainsi qu'au temps nécessaire au traitement des données. Les résultats permettront ainsi de mettre en évidence les avantages et les limites de chacun des modes étudiés et de mieux comprendre les critères à prendre en compte lors du choix d'un mode d'opération d'AES.

\section{Généralités sur la cryptographie symétrique}


\subsection{Principe de la cryptographie symétrique}

La cryptographie symétrique est une technique de chiffrement dans laquelle une même clé secrète est utilisée pour chiffrer et déchiffrer les données. L'émetteur transforme ainsi le message original, appelé texte clair, en un message chiffré qui ne peut être compris sans disposer de la clé appropriée.

Le principe général peut être représenté par la relation suivante :

\begin{equation}
C = E_K(M)
\end{equation}

où $M$ représente le message en clair, $K$ la clé secrète, $E$ la fonction de chiffrement et $C$ le message chiffré.

Le processus inverse permet de retrouver le message original :

\begin{equation}
M = D_K(C)
\end{equation}

où $D$ représente la fonction de déchiffrement.

La sécurité d'un système de chiffrement symétrique repose principalement sur la confidentialité de la clé. En effet, une personne disposant de cette clé peut généralement déchiffrer les données protégées. La gestion et la distribution des clés constituent donc des éléments essentiels dans la conception d'un système cryptographique.

Les algorithmes symétriques sont particulièrement adaptés au traitement de grandes quantités de données en raison de leur efficacité et de leur rapidité. Ils sont ainsi largement utilisés dans les systèmes informatiques, les réseaux et les protocoles de communication sécurisés.


\subsection{Chiffrement et déchiffrement}


Dans un système de cryptographie symétrique, le chiffrement et le déchiffrement constituent deux opérations complémentaires reposant sur une même clé secrète. Le chiffrement permet de transformer un message en clair en une représentation chiffrée, tandis que le déchiffrement permet de reconstruire le message original à partir du texte chiffré.

Soit $M$ un message en clair et $K$ une clé secrète. Le chiffrement de $M$ à l'aide de la clé $K$ peut être représenté par :

\begin{equation}
C = E_K(M)
\end{equation}

où $E_K$ désigne la fonction de chiffrement associée à la clé $K$ et $C$ le texte chiffré obtenu.

Le déchiffrement consiste alors à appliquer la fonction inverse afin de retrouver le message original :

\begin{equation}
M = D_K(C)
\end{equation}

où $D_K$ représente la fonction de déchiffrement.

Pour qu'un système cryptographique soit correctement fonctionnel, les opérations de chiffrement et de déchiffrement doivent satisfaire la propriété :

\begin{equation}
D_K(E_K(M)) = M
\end{equation}

Cette relation signifie que le déchiffrement d'un message précédemment chiffré avec la même clé permet de retrouver exactement le message initial.

Dans la pratique, la sécurité d'un algorithme symétrique ne doit pas reposer sur le secret de l'algorithme lui-même, mais principalement sur la protection de la clé cryptographique. Une clé compromise peut en effet permettre à un adversaire disposant du texte chiffré de retrouver les données protégées.

La cryptographie symétrique présente également l'avantage d'être relativement rapide et adaptée au chiffrement de volumes importants de données. C'est notamment pour cette raison que des algorithmes tels que l'AES sont largement utilisés dans les systèmes de communication et de stockage sécurisés.

Cependant, lorsqu'un chiffrement par blocs est utilisé pour protéger des données dont la longueur dépasse celle d'un bloc, une question supplémentaire apparaît : \textit{comment appliquer correctement l'algorithme de chiffrement à l'ensemble du message ?}

Cette problématique conduit à l'utilisation des \textbf{modes d'opération des chiffrements par blocs}, qui définissent la manière dont les différents blocs de données sont traités et combinés.


\subsection{Gestion et distribution des clés secrètes}

La sécurité d'un système de chiffrement symétrique dépend fortement de la protection de la clé secrète. En effet, contrairement au message chiffré qui peut être transmis sur un canal non sécurisé, la clé doit rester confidentielle et ne doit être accessible qu'aux entités autorisées.

Dans une communication entre deux parties, une première difficulté consiste donc à établir une clé secrète commune avant l'échange des données. Cette opération peut devenir complexe lorsque le nombre d'utilisateurs augmente. Si chaque paire d'utilisateurs doit disposer d'une clé distincte, le nombre de clés nécessaires pour $n$ utilisateurs est donné par :

\begin{equation}
N = \frac{n(n-1)}{2}
\end{equation}

Cette croissance quadratique rend la gestion des clés de plus en plus difficile dans les systèmes comportant un grand nombre d'utilisateurs. À cela s'ajoutent les problèmes liés au stockage sécurisé, au renouvellement périodique des clés, à leur distribution et à leur révocation lorsqu'elles sont compromises.

La distribution des clés constitue ainsi l'une des principales limites pratiques de la cryptographie symétrique. Plusieurs mécanismes peuvent être utilisés pour résoudre ce problème, notamment l'utilisation d'un serveur de gestion de clés ou le recours à des mécanismes de cryptographie asymétrique permettant d'établir une clé secrète de manière sécurisée.

Une fois la clé établie, le système symétrique présente toutefois un avantage important : les opérations de chiffrement et de déchiffrement sont généralement rapides et peuvent être appliquées efficacement à de grandes quantités de données.

La gestion de la clé ne constitue cependant pas la seule considération. Pour les chiffrements par blocs, il est également nécessaire de définir la manière dont les blocs successifs du message sont traités. Cette fonction est assurée par les modes d'opération, dont dépendent notamment certaines propriétés de sécurité et les performances du chiffrement.

\subsection{Modes d'opération des chiffrements par blocs}

Un chiffrement par blocs traite les données sous forme de blocs de taille fixe. Lorsqu'un message est plus long qu'un seul bloc, il est nécessaire de définir une méthode permettant d'appliquer l'algorithme de chiffrement à l'ensemble des données. Cette méthode est appelée \textit{mode d'opération}.

Un mode d'opération définit notamment la manière dont les blocs de données successifs sont traités et, selon le mode considéré, comment ils sont liés les uns aux autres. Il peut également introduire des paramètres supplémentaires tels qu'un vecteur d'initialisation (\textit{Initialization Vector}, IV) ou un nonce.

Le choix du mode d'opération ne constitue pas un simple choix d'implémentation. Il influence directement les propriétés de sécurité du système ainsi que son comportement lors du traitement des données. Certains modes assurent principalement la confidentialité, tandis que d'autres permettent également de garantir l'intégrité et l'authenticité des données.

Parmi les différents modes d'opération utilisés avec les chiffrements par blocs, on peut notamment citer ECB (\textit{Electronic Codebook}), CBC (\textit{Cipher Block Chaining}), CFB (\textit{Cipher Feedback}), OFB (\textit{Output Feedback}) et CTR (\textit{Counter Mode}). Ces modes présentent des mécanismes de fonctionnement et des propriétés de sécurité différentes.

Dans le cadre de cette étude, l'intérêt porte particulièrement sur deux modes associés à l'algorithme AES : CBC et GCM. Le mode CBC repose sur un mécanisme de chaînage entre les blocs et assure principalement la confidentialité des données. Le mode GCM, quant à lui, combine le chiffrement avec un mécanisme d'authentification permettant également de détecter certaines modifications des données.

L'étude de ces modes nécessite donc de considérer à la fois leur fonctionnement interne, leurs exigences de sécurité et leur influence sur les performances. Ces éléments constituent la base de la comparaison qui sera développée dans les sections consacrées à CBC et GCM.

\section{L'algorithme AES}

\subsection{Présentation générale d'AES}

L'\textit{Advanced Encryption Standard} (AES) est un standard de
chiffrement symétrique par blocs destiné notamment à remplacer le
\textit{Data Encryption Standard} (DES). Le développement de l'AES a
été initié par le \textit{National Institute of Standards and
Technology} (NIST) en janvier 1997, avant son adoption comme standard
en novembre 2001.

L'AES est fondé sur l'algorithme \textit{Rijndael}, conçu par les
cryptographes Joan Daemen et Vincent Rijmen. Son développement s'est
appuyé sur un processus ouvert de consultation et d'analyse auquel ont
participé des experts internationaux.

Sur le plan structurel, AES est un chiffrement par blocs itératif.
Contrairement au DES, il ne repose pas sur un réseau de Feistel.
Sa construction s'appuie principalement sur une succession de
substitutions et de transpositions, complétées par des opérations
algébriques sur les éléments manipulés par l'algorithme.

L'algorithme original \textit{Rijndael} permet de considérer plusieurs
combinaisons de longueurs de blocs et de clés, notamment 128, 192 et
256 bits. Le nombre de tours de chiffrement dépend de la configuration
retenue et peut être de 10, 12 ou 14 tours. Cette combinaison des
longueurs de bloc et de clé conduit, pour Rijndael, à plusieurs
configurations possibles.

Dans le cas de l'AES standardisé par le NIST, la taille du bloc est
fixée à 128 bits. Les trois longueurs de clé retenues sont 128, 192 et
256 bits, associées respectivement à 10, 12 et 14 tours de
transformation.

Ainsi, la structure d'AES repose sur l'application répétée de
transformations cryptographiques destinées à transformer l'état
interne du bloc au cours des différentes rondes. L'étude de ces
transformations permet de mieux comprendre le fonctionnement interne
de l'algorithme et prépare l'analyse des modes d'opération CBC et GCM
étudiés dans la suite de ce travail.

\subsection{Structure et fonctionnement d'AES}


L'algorithme AES opère sur des blocs de données en clair de 128 bits
qu'il transforme en blocs chiffrés de 128 bits à partir d'une clé
secrète de 128, 192 ou 256 bits. Le nombre de tours (\textit{rounds})
effectués au cours du chiffrement dépend de la longueur de la clé :
10 tours pour une clé de 128 bits, 12 tours pour une clé de 192 bits
et 14 tours pour une clé de 256 bits.

Le traitement interne d'AES repose sur des opérations effectuées sur
des octets et sur des mots de 32 bits. Ces opérations sont définies
à partir de structures algébriques permettant de réaliser les
différentes transformations du chiffrement.

\subsubsection{Représentation des octets dans le corps fini
$\mathrm{GF}(2^8)$}

Un octet est constitué de huit bits. Il peut être représenté par un
polynôme de degré inférieur ou égal à 7 à coefficients dans
$\mathrm{GF}(2)$ :

\begin{equation}
b(x)=b_7x^7+b_6x^6+\cdots+b_1x+b_0
\end{equation}

où chaque coefficient $b_i$ appartient à $\{0,1\}$.

Dans AES, les opérations sur les octets sont réalisées dans le corps
fini $\mathrm{GF}(2^8)$. L'addition de deux éléments correspond à
l'addition des coefficients modulo 2 et est donc équivalente à une
opération XOR bit à bit :

\begin{equation}
a(x)\oplus b(x)
\end{equation}

La multiplication est effectuée comme une multiplication
polynomiale classique, suivie d'une réduction modulo un polynôme
irréductible de degré 8. Dans AES, le polynôme utilisé est :

\begin{equation}
m(x)=x^8+x^4+x^3+x+1
\end{equation}

Ainsi, la multiplication de deux éléments $a(x)$ et $b(x)$ est
réalisée selon :

\begin{equation}
c(x)=a(x)b(x)\bmod m(x)
\end{equation}

Cette construction permet de disposer de $2^8=256$ éléments et
fournit le cadre algébrique utilisé par plusieurs transformations
internes d'AES.

\subsubsection{Représentation des mots de 32 bits}

AES manipule également des mots de 32 bits, chacun étant constitué de
quatre octets. Un tel mot peut être représenté comme un polynôme de
degré inférieur ou égal à 3 dont les coefficients appartiennent au
corps $\mathrm{GF}(2^8)$ :

\begin{equation}
B(x)=b_3x^3+b_2x^2+b_1x+b_0
\end{equation}

Les opérations sur ces mots sont alors définies dans l'anneau de
polynômes $\mathrm{GF}(2^8)[x]$. Lorsqu'une multiplication produit un
polynôme de degré supérieur à 3, celui-ci est réduit modulo :

\begin{equation}
M(x)=x^4+1
\end{equation}

Cette représentation algébrique intervient notamment dans la
transformation \textit{MixColumns}, où les quatre octets d'une colonne
de l'état sont traités comme un mot de 32 bits.

La compréhension de ces structures permet ainsi d'aborder les
principales transformations qui constituent les différentes rondes
du chiffrement AES : \textit{SubBytes}, \textit{ShiftRows},
\textit{MixColumns} et \textit{AddRoundKey}.

\subsection{Les transformations d'AES}
% ============================================================
% 2.3 LES TRANSFORMATIONS D'AES
% ============================================================

\section{L'algorithme AES}

\subsection{Les transformations d'AES}

L'algorithme AES repose sur quatre transformations fondamentales :
\textit{SubBytes}, \textit{ShiftRows}, \textit{MixColumns} et
\textit{AddRoundKey}. Ces transformations sont appliquées au tableau
d'état au cours des différentes rondes du chiffrement.

Pour AES, la taille du bloc est fixée à 128 bits, soit 16 octets.
Le tableau d'état est donc constitué de quatre lignes et de quatre
colonnes. Chaque élément du tableau représente un octet appartenant
au corps fini $\mathrm{GF}(2^8)$.

\[
State=
\begin{bmatrix}
s_{0,0}&s_{0,1}&s_{0,2}&s_{0,3}\\
s_{1,0}&s_{1,1}&s_{1,2}&s_{1,3}\\
s_{2,0}&s_{2,1}&s_{2,2}&s_{2,3}\\
s_{3,0}&s_{3,1}&s_{3,2}&s_{3,3}
\end{bmatrix}
\]

Pour AES-128, le chiffrement comporte dix rondes. Une transformation
\textit{AddRoundKey} est d'abord appliquée avant la première ronde.
Les rondes 1 à 9 utilisent les quatre transformations, tandis que
la dixième ronde ne comporte pas \textit{MixColumns}.

\begin{center}
\[
\boxed{
\begin{array}{c}
\text{Clé de chiffrement}\\
\downarrow\\
\text{Expansion de la clé}\\
\downarrow\\
K_0,K_1,\ldots,K_{10}\\[0.2cm]
\downarrow\\
\text{AddRoundKey initial}\\
\downarrow\\
\text{Rondes 1 à 9}\\
\text{SubBytes}
\rightarrow
\text{ShiftRows}
\rightarrow
\text{MixColumns}
\rightarrow
\text{AddRoundKey}\\
\downarrow\\
\text{Ronde 10}\\
\text{SubBytes}
\rightarrow
\text{ShiftRows}
\rightarrow
\text{AddRoundKey}\\
\downarrow\\
\text{Texte chiffré}
\end{array}}
\]
\end{center}


% ============================================================
% SUBBYTES
% ============================================================

\subsubsection{La transformation SubBytes}

La transformation \textit{SubBytes}, appelée également
\textit{ByteSub}, consiste à substituer indépendamment chacun des
octets du tableau d'état par son image dans une table de substitution
appelée S-Box.

Si $s_{i,j}$ désigne un octet de l'état, alors :

\[
\boxed{
s'_{i,j}=S(s_{i,j})
}
\]

où $S$ représente la S-Box d'AES.

La transformation est appliquée indépendamment à chacun des
16 octets de l'état.

\paragraph{Construction mathématique de la S-Box.}

La S-Box AES est construite à partir des propriétés du corps fini
$\mathrm{GF}(2^8)$. Pour un élément non nul $a$, on détermine d'abord
son inverse multiplicatif :

\[
a^{-1}\in\mathrm{GF}(2^8).
\]

Une transformation affine est ensuite appliquée à cet inverse.

Le principe peut être représenté par :

\[
\boxed{
a
\longrightarrow
a^{-1}
\longrightarrow
\text{transformation affine}
\longrightarrow
S(a)
}
\]

avec la convention :

\[
S(0)=0.
\]

La S-Box ainsi obtenue est bijective. Son inverse peut donc être
utilisé lors du déchiffrement :

\[
\boxed{
S^{-1}(S(a))=a
}
\]


% ------------------------------------------------------------
% TABLEAU S-BOX
% ------------------------------------------------------------

\paragraph{Table de substitution S-Box.}

La S-Box est organisée sous la forme d'un tableau $16\times16$.
Pour un octet représenté sous la forme hexadécimale
$\texttt{xy}$, le premier chiffre hexadécimal détermine la ligne
et le second détermine la colonne.

\begin{table}[H]
\centering
\caption{S-Box AES}
\renewcommand{\arraystretch}{1.15}
\setlength{\tabcolsep}{4.5pt}

\begin{tabular}{c|cccccccccccccccc}
 &0&1&2&3&4&5&6&7&8&9&A&B&C&D&E&F\\
\hline
0&63&7c&77&7b&f2&6b&6f&c5&30&01&67&2b&fe&d7&ab&76\\
1&ca&82&c9&7d&fa&59&47&f0&ad&d4&a2&af&9c&a4&72&c0\\
2&b7&fd&93&26&36&3f&f7&cc&34&a5&e5&f1&71&d8&31&15\\
3&04&c7&23&c3&18&96&05&9a&07&12&80&e2&eb&27&b2&75\\
4&09&83&2c&1a&1b&6e&5a&a0&52&3b&d6&b3&29&e3&2f&84\\
5&53&d1&00&ed&20&fc&b1&5b&6a&cb&be&39&4a&4c&58&cf\\
6&d0&ef&aa&fb&43&4d&33&85&45&f9&02&7f&50&3c&9f&a8\\
7&51&a3&40&8f&92&9d&38&f5&bc&b6&da&21&10&ff&f3&d2\\
8&cd&0c&13&ec&5f&97&44&17&c4&a7&7e&3d&64&5d&19&73\\
9&60&81&4f&dc&22&2a&90&88&46&ee&b8&14&de&5e&0b&db\\
A&e0&32&3a&0a&49&06&24&5c&c2&d3&ac&62&91&95&e4&79\\
B&e7&c8&37&6d&8d&d5&4e&a9&6c&56&f4&ea&65&7a&ae&08\\
C&ba&78&25&2e&1c&a6&b4&c6&e8&dd&74&1f&4b&bd&8b&8a\\
D&70&3e&b5&66&48&03&f6&0e&61&35&57&b9&86&c1&1d&9e\\
E&e1&f8&98&11&69&d9&8e&94&9b&1e&87&e9&ce&55&28&df\\
F&8c&a1&89&0d&bf&e6&42&68&41&99&2d&0f&b0&54&bb&16
\end{tabular}
\end{table}


% ------------------------------------------------------------
% EXEMPLE SUBBYTES
% ------------------------------------------------------------

\paragraph{Exemple de substitution d'un octet.}

Considérons l'octet :

\[
a=\texttt{53}.
\]

Le premier chiffre hexadécimal est $5$ et le second est $3$.
Nous sélectionnons donc la ligne $5$ et la colonne $3$ de la S-Box.

On obtient :

\[
\boxed{
S(\texttt{53})=\texttt{ED}
}
\]

Ainsi :

\[
\boxed{
\texttt{53}
\xrightarrow{\text{S-Box}}
\texttt{ED}
}
\]

Un autre exemple est :

\[
S(\texttt{19})=\texttt{D4}.
\]


% ------------------------------------------------------------
% EXEMPLE SUBBYTES MATRICE
% ------------------------------------------------------------

\paragraph{Application à un tableau d'état.}

Considérons le tableau :

\[
State=
\begin{bmatrix}
19&a0&9a&e9\\
3d&f4&c6&f8\\
e3&e2&8d&48\\
be&2b&2a&08
\end{bmatrix}.
\]

Chaque octet est remplacé indépendamment.

Par exemple :

\[
S(\texttt{19})=\texttt{D4},
\qquad
S(\texttt{A0})=\texttt{E0},
\]

\[
S(\texttt{9A})=\texttt{B8},
\qquad
S(\texttt{E9})=\texttt{1E}.
\]

Après application de la S-Box à tous les octets :

\[
\boxed{
\begin{bmatrix}
19&a0&9a&e9\\
3d&f4&c6&f8\\
e3&e2&8d&48\\
be&2b&2a&08
\end{bmatrix}
\overset{\text{SubBytes}}{\longrightarrow}
\begin{bmatrix}
d4&e0&b8&1e\\
27&bf&b4&41\\
11&98&5d&52\\
ae&f1&e5&30
\end{bmatrix}
}
\]


% ============================================================
% SHIFTROWS
% ============================================================

\subsubsection{La transformation ShiftRows}

La transformation \textit{ShiftRows} réalise un décalage cyclique
vers la gauche des lignes du tableau d'état.

Pour AES, les décalages sont :

\begin{table}[H]
\centering
\caption{Décalage des lignes lors de ShiftRows}
\begin{tabular}{|c|c|}
\hline
\textbf{Ligne} & \textbf{Décalage cyclique}\\
\hline
0 & 0\\
\hline
1 & 1\\
\hline
2 & 2\\
\hline
3 & 3\\
\hline
\end{tabular}
\end{table}

Pour un état général :

\[
State=
\begin{bmatrix}
a_0&a_1&a_2&a_3\\
b_0&b_1&b_2&b_3\\
c_0&c_1&c_2&c_3\\
d_0&d_1&d_2&d_3
\end{bmatrix}
\]

on obtient après \textit{ShiftRows} :

\[
\boxed{
State'=
\begin{bmatrix}
a_0&a_1&a_2&a_3\\
b_1&b_2&b_3&b_0\\
c_2&c_3&c_0&c_1\\
d_3&d_0&d_1&d_2
\end{bmatrix}
}
\]

Pour AES-128, on peut également écrire :

\[
\boxed{
s'_{i,j}=s_{i,(j+i)\bmod4}
}
\]

où $i$ représente le numéro de ligne et $j$ le numéro de colonne.


% ------------------------------------------------------------
% EXEMPLE SHIFTROWS
% ------------------------------------------------------------

\paragraph{Exemple numérique.}

Partons de l'état obtenu après \textit{SubBytes} :

\[
State=
\begin{bmatrix}
d4&e0&b8&1e\\
27&bf&b4&41\\
11&98&5d&52\\
ae&f1&e5&30
\end{bmatrix}.
\]

Les quatre lignes sont transformées comme suit :

\[
(d4,e0,b8,1e)
\longrightarrow
(d4,e0,b8,1e)
\]

\[
(27,bf,b4,41)
\longrightarrow
(bf,b4,41,27)
\]

\[
(11,98,5d,52)
\longrightarrow
(5d,52,11,98)
\]

\[
(ae,f1,e5,30)
\longrightarrow
(30,ae,f1,e5)
\]

On obtient finalement :

\[
\boxed{
\begin{bmatrix}
d4&e0&b8&1e\\
27&bf&b4&41\\
11&98&5d&52\\
ae&f1&e5&30
\end{bmatrix}
\overset{\text{ShiftRows}}{\longrightarrow}
\begin{bmatrix}
d4&e0&b8&1e\\
bf&b4&41&27\\
5d&52&11&98\\
30&ae&f1&e5
\end{bmatrix}
}
\]

Cette transformation ne modifie donc pas la valeur des octets.
Elle modifie uniquement leur position.


% ============================================================
% MIXCOLUMNS
% ============================================================

\subsubsection{La transformation MixColumns}

La transformation \textit{MixColumns} réalise un mélange des octets
au niveau de chaque colonne du tableau d'état.

Une colonne est considérée comme un vecteur de quatre éléments de
$\mathrm{GF}(2^8)$ :

\[
C=
\begin{bmatrix}
a\\
b\\
c\\
d
\end{bmatrix}.
\]

AES utilise la matrice constante :

\[
M=
\begin{bmatrix}
02&03&01&01\\
01&02&03&01\\
01&01&02&03\\
03&01&01&02
\end{bmatrix}.
\]

Le calcul est effectué dans $\mathrm{GF}(2^8)$ :

\[
\boxed{
C'=M\times C
}
\]

soit :

\[
\begin{bmatrix}
a'\\
b'\\
c'\\
d'
\end{bmatrix}
=
\begin{bmatrix}
02&03&01&01\\
01&02&03&01\\
01&01&02&03\\
03&01&01&02
\end{bmatrix}
\begin{bmatrix}
a\\
b\\
c\\
d
\end{bmatrix}.
\]

On obtient :

\[
a'=(02\otimes a)
\oplus(03\otimes b)
\oplus c
\oplus d
\]

\[
b'=a
\oplus(02\otimes b)
\oplus(03\otimes c)
\oplus d
\]

\[
c'=a
\oplus b
\oplus(02\otimes c)
\oplus(03\otimes d)
\]

\[
d'=(03\otimes a)
\oplus b
\oplus c
\oplus(02\otimes d).
\]

Le symbole $\oplus$ représente l'addition dans
$\mathrm{GF}(2^8)$, correspondant à un XOR bit à bit, tandis que
$\otimes$ représente la multiplication dans ce corps.


% ------------------------------------------------------------
% GF(2^8)
% ------------------------------------------------------------

\paragraph{Calcul dans $\mathrm{GF}(2^8)$.}

Les opérations de \textit{MixColumns} sont réalisées dans le corps
fini $\mathrm{GF}(2^8)$ à l'aide du polynôme irréductible :

\[
\boxed{
m(x)=x^8+x^4+x^3+x+1
}
\]

Un octet est représenté par un polynôme à coefficients dans
$\mathrm{GF}(2)$ et de degré inférieur ou égal à 7.

Par exemple :

\[
\texttt{57}=00111001_2
\]

correspond au polynôme :

\[
x^5+x^4+x^3+1.
\]

La multiplication dans $\mathrm{GF}(2^8)$ est effectuée par
multiplication polynomiale suivie d'une réduction modulo $m(x)$.


% ------------------------------------------------------------
% MULTIPLICATION PAR 02
% ------------------------------------------------------------

\paragraph{Exemple : multiplication par \texttt{02}.}

Considérons :

\[
a=\texttt{57}.
\]

En binaire :

\[
\texttt{57}=00111001_2.
\]

Le bit de poids fort étant nul, on effectue un décalage d'un bit
vers la gauche :

\[
01010111_2
\longrightarrow
10101110_2.
\]

On obtient :

\[
\boxed{
\texttt{02}\otimes\texttt{57}=\texttt{AE}
}
\]


% ------------------------------------------------------------
% MULTIPLICATION PAR 03
% ------------------------------------------------------------

\paragraph{Exemple : multiplication par \texttt{03}.}

Dans $\mathrm{GF}(2^8)$ :

\[
\texttt{03}\otimes a
=
(\texttt{02}\otimes a)\oplus a.
\]

Ainsi :

\[
\texttt{03}\otimes\texttt{57}
=
\texttt{AE}\oplus\texttt{57}.
\]

En binaire :

\[
\begin{array}{c}
10101110\\
\oplus\ 01010111\\
\hline
11111001
\end{array}
\]

donc :

\[
\boxed{
\texttt{03}\otimes\texttt{57}=\texttt{F9}
}
\]


% ------------------------------------------------------------
% EXEMPLE COMPLET MIXCOLUMNS
% ------------------------------------------------------------

\paragraph{Exemple complet de MixColumns.}

Considérons la colonne :

\[
C=
\begin{bmatrix}
d4\\
bf\\
5d\\
30
\end{bmatrix}.
\]

La première composante du résultat est :

\[
a'=
(02\otimes d4)
\oplus
(03\otimes bf)
\oplus
5d
\oplus
30.
\]

Les opérations dans $\mathrm{GF}(2^8)$ donnent :

\[
02\otimes d4=\texttt{B3}
\]

et :

\[
02\otimes bf=\texttt{65}.
\]

Par conséquent :

\[
03\otimes bf
=
(02\otimes bf)\oplus bf
\]

\[
=\texttt{65}\oplus\texttt{BF}
\]

\[
=\texttt{DA}.
\]

On obtient alors :

\[
a'=
\texttt{B3}
\oplus
\texttt{DA}
\oplus
\texttt{5D}
\oplus
\texttt{30}.
\]

Calculons successivement :

\[
\texttt{B3}\oplus\texttt{DA}
=
\texttt{69}
\]

\[
\texttt{69}\oplus\texttt{5D}
=
\texttt{34}
\]

\[
\texttt{34}\oplus\texttt{30}
=
\texttt{04}.
\]

Donc :

\[
\boxed{
a'=\texttt{04}
}
\]

En appliquant les mêmes opérations aux trois autres composantes :

\[
\boxed{
\begin{bmatrix}
d4\\
bf\\
5d\\
30
\end{bmatrix}
\overset{\text{MixColumns}}{\longrightarrow}
\begin{bmatrix}
04\\
66\\
81\\
e5
\end{bmatrix}
}
\]

Ainsi, chaque octet de la nouvelle colonne dépend des quatre octets
de la colonne initiale.


% ------------------------------------------------------------
% FORME POLYNOMIALE
% ------------------------------------------------------------

\paragraph{Interprétation polynomiale.}

Une colonne :

\[
(a_0,a_1,a_2,a_3)
\]

peut être représentée par le polynôme :

\[
a(x)=a_0x^3+a_1x^2+a_2x+a_3.
\]

Le polynôme constant utilisé par AES est :

\[
c(x)=
\texttt{03}x^3
\oplus
\texttt{01}x^2
\oplus
\texttt{01}x
\oplus
\texttt{02}.
\]

La transformation peut alors s'écrire :

\[
\boxed{
b(x)=c(x)a(x)\pmod{x^4+1}
}
\]

où les coefficients appartiennent à $\mathrm{GF}(2^8)$.

Cette formulation polynomiale est équivalente à la formulation
matricielle présentée précédemment.


% ============================================================
% ADDROUNDKEY
% ============================================================

\subsubsection{La transformation AddRoundKey}

La transformation \textit{AddRoundKey} consiste à combiner le tableau
d'état avec la sous-clé correspondant au round considéré.

L'opération utilisée est le XOR :

\[
\boxed{
State'=State\oplus K_i
}
\]

où $K_i$ représente la sous-clé du round $i$.

Cette opération est réalisée élément par élément entre le tableau
d'état et la sous-clé correspondante.


% ------------------------------------------------------------
% EXEMPLE XOR
% ------------------------------------------------------------

\paragraph{Exemple sur un octet.}

Considérons :

\[
State=\texttt{3A}
\]

et :

\[
K_i=\texttt{5C}.
\]

En représentation binaire :

\[
\texttt{3A}=00111010
\]

\[
\texttt{5C}=01011100.
\]

Le XOR donne :

\[
\boxed{
\begin{array}{c@{\;}cccccccc}
 & 0 & 0 & 1 & 1 & 1 & 0 & 1 & 0 \\
\oplus & 0 & 1 & 0 & 1 & 1 & 1 & 0 & 0 \\
\hline
 & 0 & 1 & 1 & 0 & 0 & 1 & 1 & 0
\end{array}
}
\]

Ainsi :

\[
\boxed{
\texttt{3A} \oplus \texttt{5C} = \texttt{66}
}
\]

% ------------------------------------------------------------
% EXEMPLE MATRICE
% ------------------------------------------------------------

\paragraph{Exemple matriciel.}

Supposons :

\[
State=
\begin{bmatrix}
19&a0&9a&e9\\
3d&f4&c6&f8\\
e3&e2&8d&48\\
be&2b&2a&08
\end{bmatrix}
\]

et une sous-clé :

\[
K_i=
\begin{bmatrix}
a0&88&23&2a\\
fa&54&a3&6c\\
fe&2c&39&76\\
17&b1&39&05
\end{bmatrix}.
\]

L'opération est effectuée élément par élément :

\[
State'=State\oplus K_i.
\]

Par exemple :

\[
\texttt{19}\oplus\texttt{A0}
=
\texttt{B9}
\]

et :

\[
\texttt{A0}\oplus\texttt{88}
=
\texttt{28}.
\]

Le résultat est donc obtenu en appliquant le même XOR à chacun
des éléments correspondants.


% ============================================================
% ENCHAINEMENT
% ============================================================

\subsubsection{Enchaînement des quatre transformations}

Les quatre transformations précédentes interviennent de manière
ordonnée au cours d'une ronde AES.

Pour une ronde standard, l'enchaînement est :

\[
\boxed{
State_i
\xrightarrow{\text{SubBytes}}
State'_i
\xrightarrow{\text{ShiftRows}}
State''_i
\xrightarrow{\text{MixColumns}}
State'''_i
\xrightarrow{\text{AddRoundKey}}
State_{i+1}
}
\]

On peut résumer le rôle de chaque transformation par le tableau
suivant :

\begin{table}[H]
\centering
\caption{Rôle des principales transformations d'AES}
\renewcommand{\arraystretch}{1.3}

\begin{tabular}{|c|c|c|}
\hline
\textbf{Transformation}
& \textbf{Nature}
& \textbf{Fonction principale}\\
\hline
\textit{SubBytes}
& Substitution non linéaire
& Introduire la non-linéarité\\
\hline
\textit{ShiftRows}
& Permutation
& Répartir les octets\\
\hline
\textit{MixColumns}
& Transformation linéaire
& Assurer la diffusion\\
\hline
\textit{AddRoundKey}
& XOR avec la sous-clé
& Introduire la clé\\
\hline
\end{tabular}
\end{table}


% ============================================================
% SCHEMA FINAL
% ============================================================

\begin{center}

\[
\boxed{
\begin{array}{c}
\text{État d'entrée}\\[0.2cm]
\downarrow\\[0.1cm]
\text{\textit{SubBytes}}\\
\downarrow\\[0.1cm]
\text{\textit{ShiftRows}}\\
\downarrow\\[0.1cm]
\text{\textit{MixColumns}}\\
\downarrow\\[0.1cm]
\text{\textit{AddRoundKey}}\\
\downarrow\\[0.2cm]
\text{État de sortie}
\end{array}
}
\]

\end{center}

La dernière ronde d'AES constitue une exception importante :
\textit{MixColumns} n'y est pas appliquée. Ainsi, pour la dernière
ronde, l'ordre des transformations est :

\[
\boxed{
\text{SubBytes}
\rightarrow
\text{ShiftRows}
\rightarrow
\text{AddRoundKey}
}
\]

L'ensemble de ces transformations permet à AES de combiner
substitution non linéaire, permutation, diffusion et incorporation
de la clé secrète. Elles constituent ainsi le mécanisme interne du
chiffrement AES par blocs, avant l'étude des modes d'opération tels
que CBC et GCM.


\section{Le mode d'opération CBC}
% ============================================================
% CHAPITRE 3 — LE MODE D'OPÉRATION CBC
% ============================================================


Le mode CBC (\textit{Cipher Block Chaining}) est un mode d'opération
utilisé avec les chiffrements symétriques par blocs, notamment AES.
Son principe consiste à chaîner les blocs de données successifs afin
que le chiffrement d'un bloc dépende également du bloc chiffré
précédent.

Dans le cas d'AES, chaque bloc possède une longueur fixe de
128 bits, soit 16 octets. Le message en clair est donc découpé en
blocs de 128 bits avant d'être traité par l'algorithme.




\subsection{Principe de fonctionnement}

Soit un message en clair constitué de $n$ blocs :

\[
P_1,\;P_2,\;\ldots,\;P_n
\]

où chaque bloc $P_i$ possède une longueur de 128 bits.

Le mode CBC utilise également un vecteur d'initialisation,
noté $IV$, de même longueur que le bloc AES :

\[
\boxed{
|IV|=128\ \text{bits}
}
\]

Le vecteur d'initialisation est utilisé pour traiter le premier bloc.
On peut alors considérer :

\[
\boxed{
C_0=IV
}
\]

Le principe général du chiffrement CBC est représenté par :

\[
\boxed{
P_i
\oplus
C_{i-1}
\longrightarrow
E_K
\longrightarrow
C_i
}
\]

où :

\begin{itemize}
    \item $P_i$ représente le bloc de texte clair numéro $i$ ;
    \item $C_i$ représente le bloc chiffré numéro $i$ ;
    \item $C_{i-1}$ représente le bloc chiffré précédent ;
    \item $C_0=IV$ pour le premier bloc ;
    \item $K$ représente la clé secrète ;
    \item $E_K$ représente la fonction de chiffrement AES avec la clé $K$ ;
    \item $\oplus$ représente l'opération XOR.
\end{itemize}

La formule générale du chiffrement CBC est donc :

\[
\boxed{
C_i=E_K(P_i\oplus C_{i-1})
}
\]

avec :

\[
\boxed{
C_0=IV
}
\]

Ainsi, le premier bloc est calculé selon :

\[
\boxed{
C_1=E_K(P_1\oplus IV)
}
\]

Le deuxième bloc dépend ensuite du premier bloc chiffré :

\[
\boxed{
C_2=E_K(P_2\oplus C_1)
}
\]

De manière générale :

\[
\boxed{
C_i=E_K(P_i\oplus C_{i-1})
}
\qquad
\text{pour }1\leq i\leq n.
\]


% ------------------------------------------------------------
% SCHÉMA GÉNÉRAL DU CHIFFREMENT CBC
% ------------------------------------------------------------

\paragraph{Représentation du chiffrement CBC.}

Le fonctionnement peut être représenté par le schéma suivant :

\[
\begin{array}{ccccc}
IV & \longrightarrow & \boxed{\oplus} & \longrightarrow &
\boxed{E_K} \longrightarrow C_1
\\[0.3cm]
&& \uparrow && \downarrow
\\[-0.1cm]
&& P_1 &&
\\[0.8cm]
C_1 & \longrightarrow & \boxed{\oplus} & \longrightarrow &
\boxed{E_K} \longrightarrow C_2
\\[0.3cm]
&& \uparrow &&
\\[-0.1cm]
&& P_2 &&
\\[0.8cm]
C_2 & \longrightarrow & \boxed{\oplus} & \longrightarrow &
\boxed{E_K} \longrightarrow C_3
\\[0.3cm]
&& \uparrow &&
\\[-0.1cm]
&& P_3 &&
\end{array}
\]

Ce schéma montre que chaque bloc de texte clair est d'abord combiné
par XOR avec le bloc chiffré précédent avant d'être soumis au
chiffrement AES.


% ------------------------------------------------------------
% EXEMPLE NUMÉRIQUE
% ------------------------------------------------------------

\paragraph{Exemple numérique simplifié.}

Pour illustrer le principe du CBC, considérons trois blocs de données :

\[
P_1,\qquad P_2,\qquad P_3
\]

et un vecteur d'initialisation $IV$.

Le premier bloc est obtenu par :

\[
P_1\oplus IV
\]

puis chiffré avec AES :

\[
\boxed{
C_1=E_K(P_1\oplus IV)
}
\]

Pour le deuxième bloc, le résultat précédent est utilisé :

\[
P_2\oplus C_1
\]

puis :

\[
\boxed{
C_2=E_K(P_2\oplus C_1)
}
\]

Enfin, le troisième bloc est traité selon :

\[
P_3\oplus C_2
\]

et :

\[
\boxed{
C_3=E_K(P_3\oplus C_2)
}
\]

L'ensemble du processus peut donc être résumé par :

\[
\boxed{
\begin{aligned}
C_1 &= E_K(P_1\oplus IV),\\
C_2 &= E_K(P_2\oplus C_1),\\
C_3 &= E_K(P_3\oplus C_2).
\end{aligned}
}
\]


% ------------------------------------------------------------
% EXEMPLE AVEC DES VALEURS HEXADÉCIMALES
% ------------------------------------------------------------

\paragraph{Illustration avec des valeurs hexadécimales.}

Considérons, à titre d'illustration du mécanisme XOR, les valeurs
suivantes :

\[
P_1=\texttt{00112233445566778899AABBCCDDEEFF}
\]

et :

\[
IV=\texttt{000102030405060708090A0B0C0D0E0F}.
\]

Avant le chiffrement AES, les deux valeurs sont combinées par XOR :

\[
P_1\oplus IV.
\]

Le calcul est effectué octet par octet :

\[
\begin{array}{c|cccccccc}
 &00&11&22&33&44&55&66&77\\
IV&00&01&02&03&04&05&06&07\\
\hline
\oplus&00&10&20&30&40&50&60&70
\end{array}
\]

et pour la seconde partie :

\[
\begin{array}{c|cccccccc}
P_1&88&99&AA&BB&CC&DD&EE&FF\\
IV&08&09&0A&0B&0C&0D&0E&0F\\
\hline
\oplus&80&90&A0&B0&C0&D0&E0&F0
\end{array}
\]

Ainsi :

\[
\boxed{
P_1\oplus IV
=
\texttt{00102030405060708090A0B0C0D0E0F0}
}
\]

Ce résultat constitue l'entrée du chiffrement AES :

\[
\boxed{
C_1=
E_K
\left(
\texttt{00102030405060708090A0B0C0D0E0F0}
\right)
}
\]

Le résultat exact de $C_1$ dépend naturellement de la clé $K$ utilisée.


% ------------------------------------------------------------
% PRINCIPE DU CHAÎNAGE
% ------------------------------------------------------------

\paragraph{Principe du chaînage.}

Le terme \textit{Cipher Block Chaining} signifie littéralement
« chaînage des blocs chiffrés ». Cette appellation provient du fait
que chaque nouveau bloc de texte clair est combiné avec le bloc
chiffré précédent.

La dépendance peut être représentée par :

\[
\boxed{
P_1\rightarrow C_1\rightarrow C_2\rightarrow C_3
}
\]

avec :

\[
C_1=E_K(P_1\oplus IV),
\]

\[
C_2=E_K(P_2\oplus C_1),
\]

\[
C_3=E_K(P_3\oplus C_2).
\]

Par conséquent, deux blocs de texte clair identiques ne produisent
pas nécessairement le même bloc chiffré lorsqu'ils apparaissent à des
positions différentes dans une même séquence CBC, puisque le bloc
précédent intervient dans le calcul.


% ------------------------------------------------------------
% CONDITION SUR LE IV
% ------------------------------------------------------------

\paragraph{Rôle du vecteur d'initialisation.}

Le vecteur d'initialisation intervient uniquement au début de la
chaîne :

\[
C_0=IV.
\]

Il permet d'éviter que le premier bloc soit toujours chiffré à partir
d'une même valeur initiale lorsque la même clé est utilisée.

Le $IV$ n'a pas besoin d'être secret. Toutefois, pour le chiffrement
CBC, son choix doit respecter les propriétés requises par le schéma
d'utilisation, notamment son caractère imprévisible au moment de la
génération.


% ------------------------------------------------------------
% TRANSITION VERS LE DÉCHIFFREMENT
% ------------------------------------------------------------

Le fonctionnement du déchiffrement CBC repose sur le mécanisme
inverse. Chaque bloc chiffré est d'abord soumis au déchiffrement AES,
puis le résultat est combiné par XOR avec le bloc chiffré précédent.

La formule du déchiffrement sera étudiée dans la sous-section
suivante.

% ============================================================
% 3.2 CHIFFREMENT ET DÉCHIFFREMENT
% ============================================================

\subsection{Chiffrement et déchiffrement}

Le principe général du chiffrement CBC a été présenté dans la
section précédente. Pour le détail du mécanisme de chiffrement,
des équations et de l'exemple numérique, se reporter à la
section~\ref{sec:cbc-principe}.

Le déchiffrement nécessite cependant une analyse particulière,
car les opérations sont effectuées dans un ordre différent de
celui du chiffrement.


% ------------------------------------------------------------
% DÉCHIFFREMENT CBC
% ------------------------------------------------------------

\subsubsection*{Déchiffrement CBC}

Lors du déchiffrement, chaque bloc chiffré $C_i$ est d'abord
traité par la fonction de déchiffrement AES $D_K$. Le résultat
obtenu est ensuite combiné par une opération XOR avec le bloc
chiffré précédent.

La relation générale du déchiffrement CBC est :

\[
\boxed{
P_i=D_K(C_i)\oplus C_{i-1}
}
\]

avec :

\[
\boxed{
C_0=IV
}
\]

Ainsi, pour le premier bloc :

\[
\boxed{
P_1=D_K(C_1)\oplus IV
}
\]

Pour le deuxième bloc :

\[
\boxed{
P_2=D_K(C_2)\oplus C_1
}
\]

et pour le troisième :

\[
\boxed{
P_3=D_K(C_3)\oplus C_2
}
\]

De manière générale :

\[
\boxed{
P_i=D_K(C_i)\oplus C_{i-1},
\qquad 1\leq i\leq n.
}
\]


% ------------------------------------------------------------
% SCHÉMA DU DÉCHIFFREMENT
% ------------------------------------------------------------

\paragraph{Représentation du déchiffrement.}

Le fonctionnement du déchiffrement CBC peut être représenté
comme suit :

\[
\begin{array}{ccccc}
C_1
&\longrightarrow&
\boxed{D_K}
&\longrightarrow&
\boxed{\oplus}
\longrightarrow P_1
\\
&&&&\uparrow
\\[-0.1cm]
&&&&IV
\\[0.8cm]

C_2
&\longrightarrow&
\boxed{D_K}
&\longrightarrow&
\boxed{\oplus}
\longrightarrow P_2
\\
&&&&\uparrow
\\[-0.1cm]
&&&&C_1
\\[0.8cm]

C_3
&\longrightarrow&
\boxed{D_K}
&\longrightarrow&
\boxed{\oplus}
\longrightarrow P_3
\\
&&&&\uparrow
\\[-0.1cm]
&&&&C_2
\end{array}
\]

Le schéma montre que le déchiffrement AES est effectué en premier,
puis que le résultat est combiné avec le bloc précédent.


% ------------------------------------------------------------
% JUSTIFICATION MATHÉMATIQUE
% ------------------------------------------------------------

\paragraph{Justification mathématique.}

Le fonctionnement du déchiffrement peut être vérifié directement
à partir de l'équation du chiffrement CBC :

\[
C_i=E_K(P_i\oplus C_{i-1}).
\]

En appliquant la fonction de déchiffrement $D_K$ aux deux membres,
on obtient :

\[
D_K(C_i)
=
D_K\left(E_K(P_i\oplus C_{i-1})\right).
\]

Puisque les fonctions $E_K$ et $D_K$ sont inverses :

\[
D_K(E_K(x))=x.
\]

On obtient alors :

\[
D_K(C_i)=P_i\oplus C_{i-1}.
\]

Il suffit ensuite d'effectuer un XOR avec $C_{i-1}$ :

\[
D_K(C_i)\oplus C_{i-1}
=
(P_i\oplus C_{i-1})\oplus C_{i-1}.
\]

Or, le XOR possède les propriétés :

\[
x\oplus x=0
\]

et :

\[
x\oplus0=x.
\]

Par conséquent :

\[
\begin{aligned}
D_K(C_i)\oplus C_{i-1}
&=
P_i\oplus C_{i-1}\oplus C_{i-1}\\
&=
P_i\oplus0\\
&=
P_i.
\end{aligned}
\]

Ainsi, on retrouve bien le bloc de texte clair initial :

\[
\boxed{
P_i=D_K(C_i)\oplus C_{i-1}.
}
\]


% ------------------------------------------------------------
% EXEMPLE NUMÉRIQUE DU DÉCHIFFREMENT
% ------------------------------------------------------------

\paragraph{Exemple numérique.}

Considérons un bloc chiffré $C_1$ et un vecteur
d'initialisation $IV$.

Après l'application du déchiffrement AES, supposons que
l'on obtienne :

\[
D_K(C_1)=
\texttt{00102030405060708090A0B0C0D0E0F0}
\]

et que :

\[
IV=
\texttt{000102030405060708090A0B0C0D0E0F}.
\]

Le bloc de texte clair est alors obtenu par :

\[
P_1=D_K(C_1)\oplus IV.
\]

Le calcul est effectué octet par octet :

\[
\begin{array}{c@{\;}cccccccc}
D_K(C_1)
&
00&10&20&30&40&50&60&70
\\
IV
&
00&01&02&03&04&05&06&07
\\
\hline
P_1
&
00&11&22&33&44&55&66&77
\end{array}
\]

Pour les huit octets suivants :

\[
\begin{array}{c@{\;}cccccccc}
D_K(C_1)
&
80&90&A0&B0&C0&D0&E0&F0
\\
IV
&
08&09&0A&0B&0C&0D&0E&0F
\\
\hline
P_1
&
88&99&AA&BB&CC&DD&EE&FF
\end{array}
\]

On retrouve donc :

\[
\boxed{
P_1=
\texttt{00112233445566778899AABBCCDDEEFF}
}
\]

Cet exemple montre que le XOR avec le vecteur d'initialisation
permet de retrouver le premier bloc de texte clair après le
déchiffrement AES.


% ------------------------------------------------------------
% CAS DES BLOCS SUIVANTS
% ------------------------------------------------------------

\paragraph{Cas des blocs suivants.}

Pour les blocs suivants, le vecteur d'initialisation n'est plus
utilisé. Il est remplacé par le bloc chiffré précédent.

Ainsi :

\[
\boxed{
P_2=D_K(C_2)\oplus C_1
}
\]

puis :

\[
\boxed{
P_3=D_K(C_3)\oplus C_2
}
\]

et ainsi de suite jusqu'au dernier bloc :

\[
\boxed{
P_n=D_K(C_n)\oplus C_{n-1}.
}
\]

Le mécanisme complet peut donc être résumé par :

\[
\boxed{
\begin{aligned}
P_1&=D_K(C_1)\oplus IV,\\
P_2&=D_K(C_2)\oplus C_1,\\
P_3&=D_K(C_3)\oplus C_2,\\
&\ \vdots\\
P_n&=D_K(C_n)\oplus C_{n-1}.
\end{aligned}
}
\]

Le déchiffrement CBC repose ainsi sur deux opérations
successives : le déchiffrement du bloc par AES, puis le XOR
avec le bloc chiffré précédent.




\subsection{Forces et limites du mode CBC}
% ============================================================
% 3.3 FORCES ET LIMITES DU MODE CBC
% ============================================================



Le mode CBC présente plusieurs propriétés intéressantes pour le
chiffrement des données. Cependant, son utilisation comporte également
certaines contraintes, notamment en matière d'authentification et de
traitement des erreurs.

% ------------------------------------------------------------
% Forces du mode CBC
% ------------------------------------------------------------

\paragraph{Forces du mode CBC}

\textbf{Chaînage des blocs.}

Le chiffrement CBC repose sur la relation :

\[
C_i = E_K(P_i \oplus C_{i-1}).
\]

Chaque bloc de texte clair dépend ainsi du bloc chiffré précédent.
Deux blocs de texte clair identiques ne produisent donc pas
nécessairement le même bloc chiffré.

\textbf{Vecteur d'initialisation.}

Pour le premier bloc, le vecteur d'initialisation est utilisé comme
valeur initiale du chaînage :

\[
C_0 = IV
\]

d'où :

\[
C_1 = E_K(P_1 \oplus IV).
\]

Le $IV$ permet ainsi d'introduire une valeur initiale dans le
processus de chiffrement.

% ------------------------------------------------------------
% Limites du mode CBC
% ------------------------------------------------------------

\paragraph{Limites du mode CBC}

\textbf{Absence d'authentification intégrée.}

Le mode CBC assure la confidentialité des données, mais ne fournit pas
à lui seul de mécanisme permettant de vérifier leur intégrité ou leur
authenticité.

\[
\boxed{
\text{Confidentialité} \neq \text{Authentification}
}
\]

Un mécanisme d'authentification complémentaire est donc nécessaire
lorsque la détection des modifications est requise.

\textbf{Nécessité du padding.}

AES utilise des blocs de 128 bits. Lorsque la longueur du message
n'est pas un multiple de cette taille, un mécanisme de remplissage
(\textit{padding}) doit être utilisé pour compléter le dernier bloc.

\textbf{Propagation des erreurs.}

Lors du déchiffrement, on a :

\[
P_i = D_K(C_i) \oplus C_{i-1}.
\]

Une modification de $C_i$ peut donc rendre $P_i$ incorrect et affecter
également le calcul de $P_{i+1}$ :

\[
P_{i+1} = D_K(C_{i+1}) \oplus C_i.
\]

Ainsi :

\[
\boxed{
C_i\ \text{modifié}
\Rightarrow
P_i\ \text{corrompu et}\ P_{i+1}\ \text{affecté}
}
\]

\textbf{Dépendance séquentielle.}

Le chiffrement d'un bloc dépend du bloc chiffré précédent :

\[
C_1 \rightarrow C_2 \rightarrow C_3
\rightarrow \cdots \rightarrow C_n.
\]

Les blocs ne peuvent donc pas être chiffrés indépendamment, ce qui
limite le parallélisme du chiffrement.

% ------------------------------------------------------------
% Synthèse
% ------------------------------------------------------------

\paragraph{Synthèse}

Les principales caractéristiques du mode CBC sont résumées dans le
tableau suivant.

\begin{table}[H]
\centering
\caption{Principales forces et limites du mode CBC}
\renewcommand{\arraystretch}{1.3}
\setlength{\tabcolsep}{6pt}

\begin{tabular}{|p{4.2cm}|p{9cm}|}
\hline
\textbf{Aspect} & \textbf{Caractéristique} \\
\hline

Chaînage &
Chaque bloc dépend du bloc chiffré précédent. \\
\hline

IV &
Le premier bloc utilise un vecteur d'initialisation. \\
\hline

Authentification &
CBC ne fournit pas d'authentification intégrée. \\
\hline

Padding &
Un remplissage peut être nécessaire pour compléter le dernier bloc. \\
\hline

Propagation des erreurs &
Une modification d'un bloc chiffré affecte le bloc correspondant et
peut affecter le suivant. \\
\hline

Parallélisation &
Le chiffrement présente une dépendance séquentielle entre les blocs. \\
\hline

\end{tabular}
\end{table}

En résumé, CBC assure la confidentialité des données grâce au
chaînage des blocs. Toutefois, l'absence d'authentification intégrée,
la gestion du padding et la propagation des modifications constituent
des limites importantes face à un adversaire actif.

% ============================================================
% 3.4 VULNÉRABILITÉS DU MODE CBC FACE AUX ATTAQUES ACTIVES
% ============================================================

\subsection{Vulnérabilités face aux attaques actives}

Contrairement à une attaque passive, dans laquelle l'adversaire se
contente d'observer les communications, une attaque active consiste
notamment à modifier, supprimer ou injecter des données dans le flux
échangé.

Le mode CBC présente certaines propriétés qui peuvent être exploitées
lorsqu'il est utilisé sans mécanisme d'authentification
complémentaire.

% ------------------------------------------------------------
% Modification du ciphertext
% ------------------------------------------------------------

\paragraph{Modification du ciphertext}

Lors du déchiffrement CBC, chaque bloc est obtenu selon :

\[
P_i = D_K(C_i) \oplus C_{i-1}.
\]

Si un adversaire modifie le bloc $C_{i-1}$ avant son arrivée chez le
destinataire, la valeur obtenue après déchiffrement devient :

\[
P'_i = D_K(C_i) \oplus C'_{i-1}.
\]

Ainsi, une modification du ciphertext peut provoquer une modification
contrôlée de certains bits du bloc de texte clair correspondant.

Cette propriété constitue une faiblesse importante lorsque CBC est
utilisé sans authentification.

% ------------------------------------------------------------
% Attaque par oracle de padding
% ------------------------------------------------------------

\paragraph{Attaque par oracle de padding}

Une autre vulnérabilité importante concerne les implémentations de CBC
utilisant un mécanisme de padding.

Dans certaines conditions, un attaquant peut modifier progressivement
le ciphertext et observer la réaction du système lors du
déchiffrement. Si le système distingue indirectement un padding valide
d'un padding invalide, ces informations peuvent être exploitées pour
obtenir progressivement des informations sur le texte clair.

Le principe général peut être représenté ainsi :

\[
\text{Ciphertext modifié}
\longrightarrow
\text{Déchiffrement}
\longrightarrow
\text{Vérification du padding}
\longrightarrow
\text{Réponse observable}
\]

La vulnérabilité ne provient donc pas d'une faiblesse du chiffrement
AES lui-même, mais de l'utilisation de CBC avec un mécanisme de
padding et d'une implémentation qui révèle une information exploitable
à l'attaquant.

% ------------------------------------------------------------
% Manipulation du IV
% ------------------------------------------------------------

\paragraph{Manipulation du vecteur d'initialisation}

Le premier bloc est déchiffré selon :

\[
P_1 = D_K(C_1) \oplus IV.
\]

Une modification du $IV$ entraîne donc une modification du premier
bloc déchiffré :

\[
P'_1 = D_K(C_1) \oplus IV'.
\]

La différence entre les deux valeurs est :

\[
P'_1 \oplus P_1 = IV' \oplus IV.
\]

Cette propriété montre pourquoi le traitement du vecteur
d'initialisation doit être réalisé avec précaution lorsqu'il est utilisé
avec CBC.

% ------------------------------------------------------------
% Conséquence générale
% ------------------------------------------------------------

\paragraph{Conséquence générale}

Les différentes vulnérabilités précédentes montrent que le chiffrement
CBC seul ne permet pas de garantir l'intégrité et l'authenticité des
données transmises.

\begin{table}[H]
\centering
\caption{Principales vulnérabilités du mode CBC}
\renewcommand{\arraystretch}{1.35}
\setlength{\tabcolsep}{6pt}

\begin{tabular}{|p{4.2cm}|p{8.8cm}|}
\hline
\textbf{Vulnérabilité} & \textbf{Conséquence} \\
\hline

Modification du ciphertext &
Possibilité de provoquer des modifications du texte clair lors du
déchiffrement. \\
\hline

Padding oracle &
Exploitation des réponses liées à la validation du padding pour
obtenir progressivement des informations sur le texte clair. \\
\hline

Manipulation du $IV$ &
Possibilité de modifier le premier bloc déchiffré lorsque le $IV$
peut être contrôlé ou modifié. \\
\hline

Absence d'authentification &
Le récepteur ne dispose pas, avec CBC seul, d'une garantie
cryptographique d'intégrité et d'authenticité du ciphertext. \\
\hline

\end{tabular}
\end{table}

Ainsi, l'utilisation de CBC dans un environnement soumis à des
attaques actives nécessite l'ajout d'un mécanisme d'authentification
approprié. Cette limitation constitue l'un des principaux éléments de
comparaison avec le mode GCM, qui associe chiffrement et
authentification dans une même construction cryptographique.


\section{Mode d'opération GCM}
\label{sec:gcm}

Le \textit{Galois Counter Mode} (GCM) est un mode d'opération
cryptographique conçu pour fournir simultanément deux propriétés
fondamentales de sécurité : la \textbf{confidentialité} des données
et leur \textbf{authentification}. Il est normalisé par le National
Institute of Standards and Technology (NIST) dans la publication
SP 800-38D.

Contrairement à un mode de chiffrement classique tel que CBC, qui
assure principalement la confidentialité, GCM associe le chiffrement
CTR (\textit{Counter Mode}) à une fonction d'authentification fondée
sur une opération de hachage dans le corps fini
$\mathrm{GF}(2^{128})$, appelée \textit{GHASH}.

Cette combinaison fait de GCM un mode de chiffrement authentifié,
communément désigné par l'acronyme \textbf{AEAD}
(\textit{Authenticated Encryption with Associated Data}).

Le fonctionnement général de GCM repose sur trois éléments :

\begin{enumerate}
    \item le chiffrement des données à l'aide du mode compteur ;
    \item l'authentification du texte chiffré et des données associées ;
    \item la génération d'une étiquette d'authentification
    (\textit{authentication tag}).
\end{enumerate}

Ainsi, le destinataire ne vérifie pas uniquement que le message peut
être déchiffré : il vérifie également que celui-ci n'a pas été modifié
pendant sa transmission.

\subsection{Chiffrement et authentification}
\label{subsec:gcm_chiffrement_authentification}

\subsubsection{Principe général}

GCM utilise généralement le chiffrement par blocs AES avec une taille
de bloc de $128$ bits. La clé secrète est notée $K$ et le nonce
(\textit{IV}, \textit{Initialization Vector}) est noté $IV$.

Une première valeur importante est calculée à partir du chiffrement
AES du bloc nul :

\begin{equation}
H = E_K(0^{128}),
\label{eq:gcm_h}
\end{equation}

où $E_K$ désigne le chiffrement AES avec la clé $K$.

La valeur $H$, appelée \textit{hash subkey}, est ensuite utilisée par
la fonction GHASH afin de calculer l'authentification des données.

Le fonctionnement de GCM peut donc être séparé conceptuellement en
deux mécanismes :

\begin{itemize}
    \item \textbf{CTR} pour assurer la confidentialité ;
    \item \textbf{GHASH} pour assurer l'authenticité et l'intégrité.
\end{itemize}

\subsubsection{Construction du compteur initial}

Lorsque le nonce possède une longueur de $96$ bits, ce qui constitue
le cas recommandé et le plus courant, le compteur initial est défini
par :

\begin{equation}
J_0 = IV \parallel 0^{31} \parallel 1,
\label{eq:gcm_j0}
\end{equation}

où $\parallel$ représente la concaténation.

Le compteur utilisé pour chiffrer les blocs suivants est obtenu par
incrémentation de la partie basse du compteur :

\begin{equation}
CTR_i = \operatorname{inc}_{32}(J_0)
\end{equation}

et, de manière générale,

\begin{equation}
CTR_i = \operatorname{inc}_{32}^{\,i}(J_0).
\end{equation}

\subsubsection{Chiffrement du message}

Soit le message en clair :

\begin{equation}
P = P_1 \parallel P_2 \parallel \cdots \parallel P_n.
\end{equation}

Pour chaque bloc $P_i$, GCM calcule une valeur pseudo-aléatoire à
partir du compteur :

\begin{equation}
S_i = E_K(CTR_i).
\end{equation}

Le texte chiffré est ensuite obtenu par une opération XOR :

\begin{equation}
C_i = P_i \oplus S_i.
\label{eq:gcm_ciphertext}
\end{equation}

Ainsi, le chiffrement d'un bloc ne dépend pas directement du bloc
précédent. Cette propriété permet notamment à GCM d'exploiter
efficacement le parallélisme disponible dans les processeurs modernes.

\subsubsection{Données associées (AAD)}

GCM permet également de traiter des données qui ne doivent pas être
chiffrées mais qui doivent néanmoins être authentifiées.

Ces données sont appelées \textbf{AAD}
(\textit{Additional Authenticated Data}).

On peut noter ces données $A$. Elles peuvent par exemple contenir :

\begin{itemize}
    \item un identifiant de session ;
    \item un numéro de séquence ;
    \item des informations d'en-tête ;
    \item des métadonnées d'un protocole réseau.
\end{itemize}

Les données $A$ ne sont donc pas chiffrées, mais leur modification
entraîne l'échec de la vérification du tag.

\subsubsection{Fonction GHASH}

L'authentification repose sur la fonction GHASH, qui effectue des
multiplications dans le corps fini :

\begin{equation}
\mathrm{GF}(2^{128}).
\end{equation}

Pour simplifier la présentation, soient :

\begin{itemize}
    \item $A$ : les données associées ;
    \item $C$ : le texte chiffré ;
    \item $H$ : la sous-clé d'authentification ;
    \item $X$ : la valeur produite par GHASH.
\end{itemize}

La fonction GHASH peut être représentée sous la forme :

\begin{equation}
X = \operatorname{GHASH}_H(A,C).
\label{eq:ghash_general}
\end{equation}

Dans le calcul réel, les données sont découpées en blocs de $128$ bits
et complétées si nécessaire. Une valeur finale contenant notamment
les longueurs de $A$ et de $C$ est également intégrée au calcul.

Pour une séquence de blocs $X_1,\ldots,X_m$, la construction peut être
exprimée récursivement par :

\begin{align}
Y_0 &= 0^{128},\\
Y_i &= (Y_{i-1} \oplus X_i)\times H,
\qquad 1\leq i\leq m.
\end{align}

La multiplication est effectuée dans $\mathrm{GF}(2^{128})$.

\subsubsection{Génération du tag d'authentification}

Une fois le résultat GHASH calculé, GCM produit le tag
d'authentification à partir de la valeur chiffrée du compteur initial.

Le tag complet peut être représenté par :

\begin{equation}
T =
E_K(J_0)
\oplus
\operatorname{GHASH}_H(A,C).
\label{eq:gcm_tag}
\end{equation}

Dans la pratique, le tag peut être tronqué à une longueur $t$ :

\begin{equation}
T_t =
\operatorname{MSB}_t
\left(
E_K(J_0)
\oplus
\operatorname{GHASH}_H(A,C)
\right),
\end{equation}

où $\operatorname{MSB}_t$ désigne les $t$ bits de poids fort.

Le résultat final de l'opération de chiffrement est donc constitué de :

\begin{equation}
\boxed{C \parallel T}
\end{equation}

avec :

\begin{itemize}
    \item $C$ : le texte chiffré ;
    \item $T$ : le tag d'authentification.
\end{itemize}

\subsubsection{Procédure complète de chiffrement}

Le fonctionnement d'Alice peut être résumé comme suit :

\begin{enumerate}
    \item Alice choisit un nonce $IV$ unique ;
    \item elle calcule la sous-clé
    $H=E_K(0^{128})$ ;
    \item elle construit le compteur initial $J_0$ ;
    \item elle chiffre les blocs du message avec le mode CTR ;
    \item elle applique GHASH aux données associées et au texte chiffré ;
    \item elle calcule le tag d'authentification ;
    \item elle transmet $(IV,C,T)$ à Bob.
\end{enumerate}

Le nonce n'est généralement pas considéré comme secret. Il doit
cependant respecter une contrainte essentielle : \textbf{il ne doit
jamais être réutilisé avec la même clé}.

% ------------------------------------------------------------
\subsection{Gestion du nonce}
\label{subsec:gcm_nonce}

La gestion du nonce constitue l'une des exigences de sécurité les
plus importantes de GCM.

Le nonce, souvent appelé $IV$ (\textit{Initialization Vector}), est
une valeur qui permet de rendre chaque opération de chiffrement
distincte.

Dans le cas recommandé d'un nonce de $96$ bits :

\begin{equation}
IV \in \{0,1\}^{96}.
\end{equation}

Le nonce n'a pas besoin d'être secret et peut donc être transmis avec
le texte chiffré. En revanche, pour une clé $K$ donnée, il doit être
\textbf{unique}.

\subsubsection{Pourquoi la réutilisation du nonce est-elle dangereuse ?}

Supposons qu'Alice chiffre deux messages différents $P$ et $P'$ avec
la même clé $K$ et le même nonce $IV$.

GCM produira alors le même flot de chiffrement :

\begin{equation}
S = E_K(CTR_i).
\end{equation}

On aura donc :

\begin{align}
C_i &= P_i \oplus S_i,\\
C'_i &= P'_i \oplus S_i.
\end{align}

Un attaquant peut alors calculer :

\begin{equation}
C_i \oplus C'_i
=
P_i \oplus P'_i.
\label{eq:nonce_reuse_xor}
\end{equation}

Le flot de chiffrement s'annule lors du XOR. L'attaquant obtient donc
une relation directe entre les deux messages en clair.

La réutilisation du nonce ne compromet donc pas uniquement la
confidentialité. Elle peut également compromettre la sécurité du
mécanisme d'authentification de GCM.

\subsubsection{Bonnes pratiques de gestion}

Pour éviter ce problème, plusieurs stratégies sont possibles :

\begin{itemize}
    \item utiliser un nonce aléatoire suffisamment grand ;
    \item utiliser un compteur strictement croissant ;
    \item utiliser une combinaison entre un identifiant de session
    et un compteur ;
    \item empêcher toute réinitialisation accidentelle du compteur.
\end{itemize}

Dans un système distribué, la gestion du nonce doit être conçue avec
une attention particulière afin d'éviter qu'un même nonce soit
généré simultanément par plusieurs instances utilisant la même clé.

\begin{tcolorbox}[
    colback=gray!5,
    colframe=black!60,
    title=\textbf{Règle fondamentale de GCM}
]
Pour une clé AES donnée, un nonce ne doit jamais être réutilisé.
Cette contrainte est essentielle à la sécurité de GCM.
\end{tcolorbox}


% ------------------------------------------------------------
\subsection{Forces et limites du mode GCM}
\label{subsec:gcm_forces_limites}

\subsubsection{Forces du mode GCM}

GCM présente plusieurs avantages importants qui expliquent son
adoption dans de nombreux protocoles modernes.

\paragraph{1. Confidentialité et authentification simultanées.}

GCM fournit dans une même construction :

\begin{equation}
\text{Confidentialité} + \text{Intégrité} + \text{Authenticité}.
\end{equation}

Il appartient ainsi à la famille des modes AEAD.

\paragraph{2. Exécution parallèle.}

Le mode CTR permet de calculer indépendamment les blocs de flux de
chiffrement :

\begin{equation}
E_K(CTR_1), E_K(CTR_2),\ldots,E_K(CTR_n).
\end{equation}

De même, les opérations internes de GCM peuvent être fortement
parallélisées. Cette caractéristique est particulièrement intéressante
pour les processeurs modernes.

\paragraph{3. Très bonnes performances.}

Les processeurs modernes peuvent disposer d'instructions matérielles
dédiées à AES et aux opérations nécessaires à GCM. Cela permet
d'obtenir des débits élevés avec une faible surcharge.

\paragraph{4. Authentification des données non chiffrées.}

Grâce aux AAD, GCM peut authentifier des informations qui doivent
rester lisibles :

\begin{equation}
A \longrightarrow \text{authentifié mais non chiffré}.
\end{equation}

Cette propriété est particulièrement utile dans les protocoles
réseau.

\subsubsection{Limites du mode GCM}

Malgré ses avantages, GCM présente certaines contraintes importantes.

\paragraph{1. Sensibilité à la réutilisation du nonce.}

C'est la principale contrainte opérationnelle de GCM. La réutilisation
d'un nonce avec la même clé peut provoquer une compromission grave de
la sécurité.

\paragraph{2. Gestion plus complexe des clés et des nonces.}

Une implémentation correcte doit garantir l'unicité des nonces.
Cette exigence devient particulièrement délicate dans les systèmes
distribués.

\paragraph{3. Risque associé aux implémentations incorrectes.}

Une erreur dans la génération du nonce, la gestion des compteurs ou
la vérification du tag peut annuler les garanties de sécurité du
mode.

\paragraph{4. Coût de l'authentification.}

Par rapport à un chiffrement seul, GCM effectue des opérations
supplémentaires pour calculer le tag. Cette surcharge reste cependant
faible sur les architectures disposant d'accélérations matérielles.

\subsubsection{Synthèse des propriétés}

Les principales caractéristiques de GCM peuvent être résumées dans
le tableau suivant.

\begin{table}[htbp]
\centering
\caption{Principales propriétés du mode GCM}
\label{tab:gcm_proprietes}
\begin{tabular}{p{4cm}p{8cm}}
\toprule
\textbf{Propriété} & \textbf{Description} \\
\midrule
Chiffrement & AES utilisé en mode compteur CTR \\
Authentification & Fonction GHASH dans $\mathrm{GF}(2^{128})$ \\
AAD & Données authentifiées sans être chiffrées \\
Tag & Permet de détecter les modifications du message \\
Nonce & Doit être unique pour une clé donnée \\
Parallélisation & Forte possibilité de parallélisation \\
Performance & Très élevée avec accélération matérielle \\
Principal risque & Réutilisation du nonce \\
\bottomrule
\end{tabular}
\end{table}

% ------------------------------------------------------------
\subsection{Résistance aux attaques actives}
\label{subsec:gcm_attaques_actives}

Une attaque active se distingue d'une attaque passive par le fait que
l'adversaire tente de modifier, supprimer, réordonner ou injecter des
messages dans le système.

GCM est conçu pour détecter ces modifications grâce au tag
d'authentification.

\subsubsection{Modification du texte chiffré}

Supposons qu'un attaquant modifie le texte chiffré :

\begin{equation}
C \longrightarrow C'.
\end{equation}

Le destinataire recalcule alors :

\begin{equation}
T' =
E_K(J_0)
\oplus
\operatorname{GHASH}_H(A,C').
\end{equation}

Comme :

\begin{equation}
C' \neq C,
\end{equation}

la valeur GHASH obtenue est normalement différente de celle calculée
par l'émetteur.

Le destinataire compare alors :

\begin{equation}
T' \stackrel{?}{=} T.
\end{equation}

Si les deux valeurs diffèrent, le message est rejeté.

\begin{tcolorbox}[
    colback=gray!5,
    colframe=black!60,
    title=\textbf{Principe de vérification}
]
Un message GCM ne doit être accepté par le destinataire que si le tag
d'authentification est correctement vérifié. Le texte chiffré ne doit
jamais être utilisé avant cette vérification.
\end{tcolorbox}

\subsubsection{Injection d'un faux message}

Un attaquant peut également tenter de fabriquer un nouveau couple
$(C',T')$.

Sans connaître la clé secrète $K$, il doit produire un tag cohérent
avec :

\begin{equation}
T' =
E_K(J_0)
\oplus
\operatorname{GHASH}_H(A,C').
\end{equation}

La difficulté de cette opération repose notamment sur le caractère
secret de la valeur $H$ et sur les propriétés cryptographiques de
GHASH.

La probabilité de réussite d'une tentative de forgery dépend notamment
de la longueur du tag et du nombre de tentatives effectuées par
l'adversaire.

\subsubsection{Attaque par modification des AAD}

Les données associées sont également protégées contre la modification.

Si :

\begin{equation}
A \longrightarrow A',
\end{equation}

alors le calcul :

\begin{equation}
\operatorname{GHASH}_H(A',C)
\end{equation}

produira normalement une valeur différente de :

\begin{equation}
\operatorname{GHASH}_H(A,C).
\end{equation}

Le tag ne correspondra donc plus et le message sera rejeté.

Cette propriété permet de protéger des champs importants d'un
protocole sans avoir à les chiffrer.

\subsubsection{Attaque par rejeu}

Il convient cependant de distinguer l'intégrité d'un message de la
protection contre le rejeu.

GCM permet de détecter la modification d'un message, mais le mode
GCM seul ne constitue pas nécessairement un mécanisme complet de
protection contre les attaques par rejeu.

Pour empêcher qu'un ancien message authentique soit retransmis, le
protocole doit généralement utiliser :

\begin{itemize}
    \item un numéro de séquence ;
    \item un compteur ;
    \item un identifiant unique de message ;
    \item ou une fenêtre de rejeu.
\end{itemize}

Ces informations peuvent être placées dans les AAD afin qu'elles
soient elles-mêmes authentifiées.

\subsubsection{Bilan de la résistance aux attaques actives}

Le niveau de protection fourni par GCM peut être résumé ainsi :

\begin{table}[htbp]
\centering
\caption{GCM face aux principales attaques actives}
\label{tab:gcm_attaques}
\begin{tabular}{p{5cm}p{7cm}}
\toprule
\textbf{Attaque} & \textbf{Protection fournie par GCM} \\
\midrule
Modification du ciphertext
& Détectée par la vérification du tag \\

Modification des AAD
& Détectée par GHASH et le tag \\

Injection d'un faux message
& Rendue difficile sans la clé secrète \\

Réutilisation d'un nonce
& Dangereuse ; peut compromettre la sécurité \\

Rejeu
& Nécessite un mécanisme supplémentaire au niveau du protocole \\
\bottomrule
\end{tabular}
\end{table}

En définitive, GCM offre une protection robuste contre la falsification
des messages à condition que ses exigences fondamentales soient
respectées. La sécurité du mode dépend en particulier de la
\textbf{non-réutilisation des nonces}, de la protection de la clé
secrète et de la vérification systématique du tag avant l'utilisation
du message déchiffré.



\section{Analyse comparative de CBC et GCM}
\label{sec:comparaison_cbc_gcm}

Après avoir étudié séparément le fonctionnement du mode CBC et celui
du mode GCM, il est nécessaire de comparer leurs propriétés afin de
mettre en évidence leurs différences en matière de sécurité, de
performances et de résistance aux attaques.



\subsection{Comparaison du niveau de sécurité}
\label{subsec:comparaison_securite}

% Contenu à développer dans la section suivante.


\subsection{Comparaison face aux attaques actives}
\label{subsec:comparaison_attaques}


\subsection{Comparaison des performances}
\subsection{Comparaison des performances}
\label{subsec:comparaison_performances}

% Contenu à développer.

\subsection{Avantages et limites de chaque mode}
\subsection{Avantages et limites de chaque mode}
\label{subsec:avantages_limites}

% Contenu à développer.



\section{Expérimentations et résultats}

\subsection{Environnement expérimental}
\subsection{Scénarios de test}
\subsection{Mesure des performances}
\subsection{Analyse des résultats}

\section{Conclusion et perspectives}



\end{document}
